\section*{Revision History}

This framework was developed across four rounds of external review, each targeting a different weakness in the prior draft.

\begin{center}
\begin{tabular}{|p{0.08\textwidth}|p{0.20\textwidth}|>{\small\raggedright\arraybackslash}p{0.60\textwidth}|}
\hline
Round & Focus & What changed \\
\hline
1 & Evidentiary accuracy & Corrected overgeneralized claims, misattributed sources (Kahneman \& Klein, Cynefin, the Bitter Lesson, Parasuraman/Sheridan/Wickens), and a factual error in the Terminus-4B case study. Named and dated the OpenAI/Hugging Face containment incident precisely rather than leaving it anonymous. Softened claims about privilege waiver, output reproducibility, and model-confidence calibration that had been stated more strongly than the evidence supports. \\
\hline
2 & Systems safety & Identified a genuine structural gap: the framework had substantial coverage of where errors originate and why verification is hard, but no explicit account of how an erroneous or uncertain state propagates through a socio-technical system and becomes --- or fails to become --- real harm. Added the material now in Sections 11 through 14: hazard-first analysis, propagation topology, verification as an assurance argument, and operational resilience. \\
\hline
3 & Purpose and standard of success & Observed that a long catalogue of failure modes, without a stated positive objective, reads as an argument against ever deploying anything. Added Section 2, establishing comparative reliability against a strong human baseline --- not perfection --- as the framework's actual standard of success. \\
\hline
4 & Job definition & Found that job definition, on which every later section depends, had been compressed into a single sentence. Expanded Section 3 into a full treatment distinguishing the job, the decision it informs, the task assigned, the output produced, and the effect that follows --- and introduced job drift as a named failure mode. \\
\hline
5 & Navigability & Found the report dense enough that a senior, time-pressed reader had no way to tell which sections carry the actual decision versus which build and pressure-test it. Added the Executive Guide immediately after the Abstract, mapping each of the six movements to the specific costly misdiagnosis it exists to prevent, without shortening or simplifying the underlying material. \\
\hline
\end{tabular}
\end{center}

The open questions each round surfaced are collected in Section 15, rather than scattered through the body text, so that unresolved disagreements remain visible rather than being quietly smoothed over.
